\documentclass[11pt,letterpaper]{article}
\usepackage[utf8]{inputenc}
\usepackage{geometry}
\usepackage{setspace}
\usepackage{enumitem}

\geometry{letterpaper, margin=1in}
\doublespacing

\setlength{\parindent}{0.5in}
\setlength{\parskip}{0pt}

\begin{document}

\thispagestyle{empty}

\begin{center}
    \textbf{\large CLAIMS}
\end{center}

\vspace{0.4cm}
\noindent What is claimed is:

\begin{enumerate}[label=\textbf{\arabic*.}, leftmargin=0.5in, itemsep=0.4cm]
    \item A method for negotiating a data encoding scheme between two devices connected by a parallel communication interface, the method comprising the steps of:
    \begin{itemize}[leftmargin=0.3in, noitemsep]
        \item establishing a communication session using a default, pre-defined encoding map;
        \item transmitting, from a first device to a second device, a capabilities message listing one or more available high-density encoding maps;
        \item transmitting, from the second device to the first device, a selection message indicating a chosen encoding map from the list; and
        \item conducting subsequent data transmission for the session using the chosen high-density encoding map.
    \end{itemize}

    \item The method of claim 1, wherein the negotiation occurs at the Physical Layer (Layer 1) or Data-Link Layer (Layer 2) of the OSI model.

    \item The method of claim 1, wherein the default encoding map is a low-cardinality map and the chosen encoding map is a high-cardinality map.

    \item The method of claim 1, wherein the parallel communication interface comprises a two-dimensional array of pixels.

    \item The method of claim 4, wherein the default encoding map represents a binary state using black and white pixel values.

    \item The method of claim 4, wherein the chosen high-density encoding map maps multi-byte data blocks to unique pixel color values.

    \item The method of claim 1, wherein the parallel communication interface is a massively parallel optical interface as described in U.S. Patent Application No. 19/559,705.
\end{enumerate}

\end{document}
